\section{Discussion}
\label{sec:discussion}

\subsection{Gauge-free, track-free storage}
The water-level-substitution result (Section~\ref{sec:res_ladder}) shows that
replacing the in-situ gauge with SWOT altimetry changes reconstructed storage by only
$\sim$5\% relative to field truth. Because SWOT observes by swath rather than along a
nadir track, it removes the two constraints that keep small reservoirs off nadir
altimetry's map: the need for a gauge, and the need for a satellite ground track to
actually cross the water body. Notably, DAHITI could add these small reservoirs only
once SWOT existed, so wide-swath altimetry is what makes them observable at all.
Rosamarina makes the case concretely: its gauge suffered two separate stuck episodes
at two different flat readings, spanning ${\sim}19$ months combined across
2024--2026, and only the SWOT level --- corroborated by the optical masks ---
recovered the true drawdown, without which the optical reconstruction collapsed to a
flat, unphysical surface (Section~\ref{sec:res_xsensor}). For this specific reservoir, the satellite source was
not merely as good as the gauge but strictly better; the reverse can also hold,
however, since at the low-A/P Ancipa the SWOT water level itself degrades (RMSE
$12$~m, Section~\ref{sec:res_ladder}) --- so which source is more reliable is
reservoir-dependent, not a general property of SWOT over gauges.

\subsection{What the method sees: the observability envelope}
The reconstruction is bounded by the drawdown-exposed band. A reservoir that draws
down close to its dead pool, or that is filling from empty, exposes most of its basin
and can be reconstructed almost completely; a reservoir kept near full exposes only a
thin upper band. Consequently, deep-zone sedimentation is invisible to SAR and the
band capacity change is a lower bound on total loss (Section~\ref{sec:res_change}).
The 2025--2026 Sicilian drought-to-refill event, by drawing most reservoirs down
sharply and refilling them within months, was particularly favourable to the method:
it exposed an unusually wide elevation range in a short, temporally consistent window
(Section~\ref{sec:intro}).

\subsection{Reliability scales with A/P}
The area-to-perimeter ratio that predicts SAR water-area reliability
\citep{PaperOne} also predicts bathymetric-reconstruction reliability: narrow,
low-A/P reservoirs (Olivo, Ancipa, Nicoletti) show the largest apparent capacity
losses in Section~\ref{sec:res_change} and, where an independent curve exists
(Nicoletti), the widest gap between SAR and truth (Section~\ref{sec:res_curves});
Ancipa additionally has a degraded SWOT water level (RMSE $12$~m). The same signature
appears in the cross-sensor comparison (Section~\ref{sec:res_xsensor}): the PlanetScope
and SAR reconstructions agree to $1.3$--$3.3$~m at Pozzillo, Poma and Rosamarina but
diverge ($6.1$~m) at the low-A/P Ancipa, which is also revisit-rate limited (only 9
Sentinel-1 acquisitions span its drawdown-refill window, all already used, so no
denser resampling is possible there). This is unlikely to be a classifier
artifact: every reservoir here is masked with the same per-scene VV-Otsu classifier
(Section~\ref{sec:data}), whose dual-polarisation SVM alternative performs comparably
overall \citep{PaperOne}, though we have not tested SVM masks specifically at Ancipa.
Multiple independent lines of evidence --- apparent capacity loss, curve validation
gap, degraded SWOT water level, and PlanetScope cross-sensor divergence --- therefore
converge on the same low-A/P reservoirs, so A/P offers an a-priori screen for where
the fully remote pipeline can be trusted.

\subsection{Systematic bias and its correction}
The elevation-dependent bias (Section~\ref{sec:res_garcia}) is intrinsic to the
level-slicing reconstruction. A single field calibration transfers to a similarly
sedimenting reservoir (Rosamarina) but not to an anomalous one (Poma), so we treat it
as a partially-correctable systematic with a $\pm11.3$~pp transfer uncertainty rather
than a universal correction; more field surveys would tighten it.

\subsection{Limitations and future work}
Three limitations bound the present results. First, the SWOT record begins only in mid
2023, so the fully-remote storage series is still short; this will lengthen with the
mission. Second, the validation sample is small (nine reservoirs, one echo-sounder
survey), although the cross-sensor and water-level-substitution experiments partly
compensate by adding independent evidence on the same reservoirs. Third, and most
fundamental, the deep permanently-wet zone is never exposed and so is invisible to any
waterline method; where a modern survey exists it can be added to close the AEV curve,
but in general the SAR estimate remains a lower bound on total capacity loss.

Two avenues would strengthen the validation directly. A field campaign combining an
uncrewed surface vehicle (echo-sounding the deep zone) with drone photogrammetry (the
exposed margin) would yield a complete, independently-measured basin DEM --- turning the
lower bound into a measured total and extending the bias calibration beyond the single
Garcia survey; the low-A/P reservoirs (Ancipa) and a sedimenting mid-A/P case
(Rosamarina) are the most informative targets. Separately, Sentinel-1 interferometry
over the \emph{exposed margin} could refine waterline delineation (interferometric
coherence) and probe compaction of freshly-exposed sediment (differential InSAR); we
stress that interferometry cannot retrieve open-water level, because the water surface
decorrelates between passes, so it is a complement to the margin geometry, not a
water-level source.
